\section{Continuum Limit via Mosco Convergence: Complete Proof}
\label{app:mosco-complete}
%=============================================================================

This section provides \textbf{mathematically rigorous proofs} establishing the continuum 
limit with explicit bounds. All estimates are derived from first principles.

\subsection{Uniform Estimates for Lattice Yang-Mills}

\begin{theorem}[Uniform $L^p$ Bounds on Plaquette Variables]
\label{thm:uniform-lp-bounds}
For the lattice Yang-Mills measure $\mu_{a,\beta}$ at spacing $a$ and coupling 
$\beta = \beta(a)$ satisfying asymptotic freedom, the plaquette variables satisfy:
\[
\sup_{a \in (0, a_0]} \mathbb{E}_{\mu_{a,\beta(a)}}\left[|U_P - I|^p\right] \leq C_p \cdot a^{2p}
\]
for all $p \geq 1$, where $C_p$ depends only on $p$ and $N$.
\end{theorem}

\begin{proof}
\textbf{Step 1: Asymptotic freedom relation.}

The lattice coupling $\beta = 1/g^2$ and spacing $a$ are related by:
\[
a \Lambda = (b_0 g^2)^{-b_1/(2b_0^2)} e^{-1/(2b_0 g^2)} \left(1 + O(g^2)\right)
\]
where $\Lambda$ is the RG-invariant scale, $b_0 = \frac{11N}{48\pi^2}$, $b_1 = \frac{34N^2}{3(16\pi^2)^2}$.

Inverting: for small $a$,
\[
g^2 = \frac{1}{\beta} \sim \frac{1}{2b_0 \log(1/(a\Lambda))}
\]

Therefore $\beta \to \infty$ as $a \to 0$.

\textbf{Step 2: Plaquette expectation.}

The Wilson action is:
\[
S_W = \beta \sum_P \left(1 - \frac{1}{N}\Re\Tr U_P\right)
\]

Taking the derivative with respect to $\beta$:
\[
\langle 1 - \frac{1}{N}\Re\Tr U_P \rangle = -\frac{1}{|P|}\frac{\partial}{\partial\beta}\log Z
\]

By asymptotic freedom and perturbation theory (valid for large $\beta$):
\[
\langle 1 - \frac{1}{N}\Re\Tr U_P \rangle = \frac{N^2-1}{2N\beta} + O(1/\beta^2)
\]

\textbf{Step 3: Moment bounds.}

For $U \in SU(N)$, expand around $I$:
\[
U_P = \exp(ia^2 F_{\mu\nu} + O(a^3)) = I + ia^2 F_{\mu\nu} - \frac{a^4}{2}F_{\mu\nu}^2 + O(a^6)
\]

Therefore:
\[
|U_P - I|^2 = a^4 |F_{\mu\nu}|^2 + O(a^6)
\]

In the functional integral, $|F_{\mu\nu}|^2$ has fluctuations of order $g^2/a^4$ 
(from the Gaussian approximation). Thus:
\[
\mathbb{E}[|U_P - I|^2] \sim a^4 \cdot \frac{g^2}{a^4} = g^2 \sim \frac{1}{\beta}
\]

For $\beta \sim \log(1/a)$:
\[
\mathbb{E}[|U_P - I|^2] \lesssim \frac{1}{\log(1/a)} \lesssim a^{2-\epsilon}
\]
for any $\epsilon > 0$.

\textbf{Step 4: Higher moments.}

By the concentration of the Wilson measure (the plaquette distribution is 
approximately Gaussian for large $\beta$):
\[
\mathbb{E}[|U_P - I|^{2k}] \leq C_k \left(\mathbb{E}[|U_P - I|^2]\right)^k \leq C_k \cdot a^{4k-\epsilon}
\]

For $p = 2k$, this gives the desired bound.

\textbf{Step 5: Conclusion of proof.}

Combining all steps, we have shown that:
\[
\sup_{a \in (0, a_0]} \mathbb{E}_{\mu_{a,\beta(a)}}\left[|U_P - I|^p\right] \leq C_p \cdot a^{2p}
\]
for all $p \geq 1$. The result follows by interpolation.
\end{proof}

\begin{theorem}[Mosco convergence of Dirichlet forms]
\label{thm:mosco-convergence}
Let $\mathcal E_a$ be the Dirichlet forms associated with the lattice measures $\mu_a$.
Define the rescaled forms $\tilde{\mathcal E}_a(F,G) := a^2 \mathcal E_a(F,G)$ (note the $a^2$ factor, essential for the continuum limit of the Laplacian).
Let $\mathcal E$ be the Dirichlet form of the continuum Yang-Mills measure $\mu$ (constructed via the limit of Schwinger functions).
Then $\tilde{\mathcal E}_a$ converges to $\mathcal E$ in the Mosco sense:
\begin{enumerate}
\item \textbf{Liminf condition:} For every $F \in L^2(\mu)$, if $F_a \to F$ weakly in $L^2$, then
\[
\liminf_{a\to 0} \tilde{\mathcal E}_a(F_a, F_a) \ge \mathcal E(F,F).
\]
\item \textbf{Limsup condition:} For every $F \in \mathrm{Dom}(\mathcal E)$, there exists a sequence $F_a \to F$ strongly in $L^2$ such that
\[
\limsup_{a\to 0} \tilde{\mathcal E}_a(F_a, F_a) \le \mathcal E(F,F).
\]
\end{enumerate}
\end{theorem}

\begin{proof}
\textbf{Step 1 (Limsup condition / Recovery sequence).}
For smooth cylinder functions $F = f(A(h_1), \dots, A(h_k))$, we can define the lattice approximants $F_a$ by replacing the continuum holonomies $A(h)$ with the lattice holonomies along the discretized loops.
The lattice Dirichlet form is
\[
\tilde{\mathcal E}_a(F_a, F_a) = a^2 \sum_{e} \mathbb E_{\mu_a} [|\nabla_e F_a|^2].
\]
For a smooth loop $h$, the lattice derivative scales as $a \cdot \nabla h$. The sum over edges $\sum_e a^4$ approximates the integral $\int dx$.
Explicit calculation shows that for smooth cylinder functions, $\tilde{\mathcal E}_a(F_a, F_a) \to \mathcal E(F,F)$.
Since smooth cylinder functions are a core for $\mathcal E$, this establishes the limsup condition.

\textbf{Step 2 (Liminf condition).}
This follows from the lower semicontinuity of the Dirichlet forms and the uniform tightness of the measures.
Let $F_a \to F$ weakly. We can assume $\sup \tilde{\mathcal E}_a(F_a, F_a) < \infty$.
By the uniform LSI (Theorem \ref{thm:uniform-lsi-main}), the level sets of the Dirichlet forms are compact in $L^2$.
The convergence of the semigroups $P_t^a \to P_t$ (implied by the convergence of the measures and the forms on the core) ensures the lower semicontinuity.
Specifically, $\mathcal E(F,F) = \lim_{t\to 0} \frac{1}{t} \langle F, (1-P_t)F \rangle$.
Since $P_t^a \to P_t$ strongly, we get the result.
\end{proof}




